\documentclass[10pt,a4paper]{amsart}
\usepackage{amsmath,amssymb,amsthm,graphicx}
\usepackage{verbatim}
\newenvironment{reference}{\comment}{\endcomment}
%\newenvironment{reference}{}{}
\newenvironment{slogan}{\comment}{\endcomment}
\newenvironment{history}{\comment}{\endcomment}

% For commutative diagrams we use Xy-pic
\usepackage[all]{xy}

% We use 2cell for 2-commutative diagrams.
\xyoption{2cell}
% Theorem environments.
%
\theoremstyle{plain}
\newtheorem{theorem}[subsection]{Theorem}
\newtheorem{proposition}[subsection]{Proposition}
\newtheorem{lemma}[subsection]{Lemma}

\theoremstyle{definition}
\newtheorem{definition}[subsection]{Definition}
\newtheorem{example}[subsection]{Example}
\newtheorem{exercise}[subsection]{Exercise}
\newtheorem{situation}[subsection]{Situation}

\theoremstyle{remark}
\newtheorem{remark}[subsection]{Remark}
\newtheorem{remarks}[subsection]{Remarks}

\numberwithin{equation}{subsection}

% Macros
%
\def\lim{\mathop{\mathrm{lim}}\nolimits}
\def\colim{\mathop{\mathrm{colim}}\nolimits}
\def\Spec{\mathop{\mathrm{Spec}}}
\def\Hom{\mathop{\mathrm{Hom}}\nolimits}
\def\Ext{\mathop{\mathrm{Ext}}\nolimits}
\def\SheafHom{\mathop{\mathcal{H}\!\mathit{om}}\nolimits}
\def\SheafExt{\mathop{\mathcal{E}\!\mathit{xt}}\nolimits}
\def\Sch{\mathit{Sch}}
\def\Mor{\mathop{\mathrm{Mor}}\nolimits}
\def\Ob{\mathop{\mathrm{Ob}}\nolimits}
\def\Sh{\mathop{\mathit{Sh}}\nolimits}
\def\NL{\mathop{N\!L}\nolimits}
\def\CH{\mathop{\mathrm{CH}}\nolimits}
\def\proetale{{pro\text{-}\acute{e}tale}}
\def\etale{{\acute{e}tale}}
\def\QCoh{\mathit{QCoh}}
\def\Ker{\mathop{\mathrm{Ker}}}
\def\Im{\mathop{\mathrm{Im}}}
\def\Coker{\mathop{\mathrm{Coker}}}
\def\Coim{\mathop{\mathrm{Coim}}}

% Boxtimes
%
\DeclareMathSymbol{\boxtimes}{\mathbin}{AMSa}{"02}

%
% Macros for moduli stacks/spaces
%
\def\QCohstack{\mathcal{QC}\!\mathit{oh}}
\def\Cohstack{\mathcal{C}\!\mathit{oh}}
\def\Spacesstack{\mathcal{S}\!\mathit{paces}}
\def\Quotfunctor{\mathrm{Quot}}
\def\Hilbfunctor{\mathrm{Hilb}}
\def\Curvesstack{\mathcal{C}\!\mathit{urves}}
\def\Polarizedstack{\mathcal{P}\!\mathit{olarized}}
\def\Complexesstack{\mathcal{C}\!\mathit{omplexes}}
% \Pic is the operator that assigns to X its picard group, usage \Pic(X)
% \Picardstack_{X/B} denotes the Picard stack of X over B
% \Picardfunctor_{X/B} denotes the Picard functor of X over B
\def\Pic{\mathop{\mathrm{Pic}}\nolimits}
\def\Picardstack{\mathcal{P}\!\mathit{ic}}
\def\Picardfunctor{\mathrm{Pic}}
\def\Deformationcategory{\mathcal{D}\!\mathit{ef}}
\usepackage{cleveref}
\crefname{originallemma}{original lemma}{original lemmas}
\Crefname{originallemma}{Original lemma}{Original lemmas}
\makeatletter
\newlabel{homology-lemma-first-quadrant-ss}{{0132}{source}{}{sourcepointer-1}{}}
\newlabel{homology-lemma-first-quadrant-ss@cref}{{[originallemma][0][]0132}{[1][0][]source}{}{}{}}
\makeatother
\crefname{originallemma}{original lemma}{original lemmas}
\Crefname{originallemma}{Original lemma}{Original lemmas}
\makeatletter
\newlabel{hypercovering-lemma-cech-spectral-sequence}{{01GY}{source}{}{sourcepointer-2}{}}
\newlabel{hypercovering-lemma-cech-spectral-sequence@cref}{{[originallemma][0][]01GY}{[1][0][]source}{}{}{}}
\makeatother
\crefname{originallemma}{original lemma}{original lemmas}
\Crefname{originallemma}{Original lemma}{Original lemmas}
\makeatletter
\newlabel{sites-cohomology-lemma-w-contractible}{{0946}{source}{}{sourcepointer-3}{}}
\newlabel{sites-cohomology-lemma-w-contractible@cref}{{[originallemma][0][]0946}{[1][0][]source}{}{}{}}
\makeatother
\title{Cohomology on the pro-etale site is computed by the total complex of a hypercovering with weakly contractible terms}
\author{The Stacks Project Authors}
\date{stacks-project:proetale.tex:local-snapshot-20261012; GNU FDL 1.2 or later}
\begin{document}
\maketitle
\noindent\textit{Editorial note.} All delivered mathematical text is the author's own, copied from the pinned Stacks Project chapter file proetale.tex; no statement or proof step was rewritten, reordered or removed.\par\medskip
\noindent\textit{Editorial note.} The counted proof is the author's complete argument: the single sheaf case through the Cech spectral sequence degenerated by the vanishing of higher cohomology over weakly contractible objects, and the passage to bounded below complexes by an injective resolution compared through the spectral sequences of a double complex.\par\medskip
\noindent\textit{Editorial note.} The label lemma-compute-cohomology names the counted statement, and the proof environment immediately following it in the source is delivered verbatim as its proof.\par\medskip
\noindent\textit{Editorial note.} The counted statement is carried by the author's own lemma environment in the chapter source.\par\medskip
\noindent\textit{Editorial note.} Pre-open gate: python operations/agent-authoring/tools/precard\_gate.py --source proetale.tex --label lemma-compute-cohomology --start 3282 --end 3327 --local-id 3590 returned RESULT=OK (L3590 free, occupancy FREE, public interval clear). Candidate taken from sources/stacks-proetale/candidates-halfA.json.\par\medskip
\noindent\textit{Editorial note.} The statement and proof are copied from the pinned chapter file; the cited inputs are kept as external pointers with their canonical Stacks tags.\par\medskip
\noindent\textit{Allowed standard prerequisites (editorial).}\begin{itemize}
\item Cech spectral sequence of a hypercovering. Hypercoverings, the lemma producing the spectral sequence computing the cohomology of a sheaf from its alternating sections over a hypercovering. Supplies the identification of R-Gamma of a single sheaf with the alternating complex of its sections over the hypercovering.
\item Vanishing of the higher cohomology over a weakly contractible object. Cohomology on Sites, the lemma that a weakly contractible object has vanishing higher cohomology with coefficients in any sheaf. Supplies the vanishing hypothesis that makes the spectral sequence degenerate in the case of a single sheaf.
\item Spectral sequences of a double complex of a map of complexes. Homology, the lemma giving the two spectral sequences of a first quadrant double complex. Supplies the quasi-isomorphism criterion for the map of total complexes and the identification of the total complex of the injective resolution as computing R-Gamma.
\end{itemize}\medskip
\ifdefined\hypertarget\hypertarget{sourcepointer-1}{}\fi
\noindent\textit{Source reference.} Original lemma 0132: sources/stacks-chow/tags.txt, source lines 1420--1420. Prerequisite pointer invoked at proof line(s) 3319; it names a standard fact of a sibling chapter that the counted argument uses, and every step of the counted proof built on it is retained verbatim in the delivered body.\par\medskip
\ifdefined\hypertarget\hypertarget{sourcepointer-2}{}\fi
\noindent\textit{Source reference.} Original lemma 01GY: sources/stacks-chow/tags.txt, source lines 1923--1923. Prerequisite pointer invoked at proof line(s) 3298; it names a standard fact of a sibling chapter that the counted argument uses, and every step of the counted proof built on it is retained verbatim in the delivered body.\par\medskip
\ifdefined\hypertarget\hypertarget{sourcepointer-3}{}\fi
\noindent\textit{Source reference.} Original lemma 0946: sources/stacks-chow/tags.txt, source lines 11539--11539. Prerequisite pointer invoked at proof line(s) 3304; it names a standard fact of a sibling chapter that the counted argument uses, and every step of the counted proof built on it is retained verbatim in the delivered body.\par\medskip
\begin{lemma}
\label{lemma-compute-cohomology}
Let $X$ be a scheme. Let $E \in D^+(X_\proetale)$ be represented by
a bounded below complex $\mathcal{E}^\bullet$ of abelian sheaves.
Let $K$ be a hypercovering of $U \in \Ob(X_\proetale)$ with
$K_n = \{U_n \to U\}$ where $U_n$ is a weakly contractible object of
$X_\proetale$. Then
$$
R\Gamma(U, E) = \text{Tot}(s(\mathcal{E}^\bullet(K)))
$$
in $D(\textit{Ab})$.
\end{lemma}

\begin{proof}
If $\mathcal{E}$ is an abelian
sheaf on $X_\proetale$, then the spectral sequence of
Hypercoverings, Lemma \ref{hypercovering-lemma-cech-spectral-sequence}
implies that
$$
R\Gamma(X_\proetale, \mathcal{E}) = s(\mathcal{E}(K))
$$
because the higher cohomology groups of any sheaf over $U_n$ vanish, see
Cohomology on Sites, Lemma \ref{sites-cohomology-lemma-w-contractible}.

\medskip\noindent
If $\mathcal{E}^\bullet$ is bounded below, then we can choose an injective
resolution $\mathcal{E}^\bullet \to \mathcal{I}^\bullet$ and consider
the map of complexes
$$
\text{Tot}(s(\mathcal{E}^\bullet(K)))
\longrightarrow
\text{Tot}(s(\mathcal{I}^\bullet(K)))
$$
For every $n$ the map $\mathcal{E}^\bullet(U_n) \to \mathcal{I}^\bullet(U_n)$
is a quasi-isomorphism because taking sections over $U_n$ is exact.
Hence the displayed map is a quasi-isomorphism by one of the spectral
sequences of
Homology, Lemma \ref{homology-lemma-first-quadrant-ss}.
Using the result of the first paragraph we see that for every $p$
the complex $s(\mathcal{I}^p(K))$ is acyclic in degrees $n > 0$ and
computes $\mathcal{I}^p(U)$ in degree $0$. Thus the other spectral
sequence of
Homology, Lemma \ref{homology-lemma-first-quadrant-ss}
shows $\text{Tot}(s(\mathcal{I}^\bullet(K)))$ computes
$R\Gamma(U, E) = \mathcal{I}^\bullet(U)$.
\end{proof}

\end{document}
